\documentclass[12pt]{article}
\usepackage{amsfonts,amsthm,amsmath,amssymb,mathtools}
\usepackage{mathrsfs}
\usepackage{enumitem}
\usepackage{tikz-cd}
\usepackage{geometry}
\usepackage{mlmodern}

\geometry{letterpaper, portrait, margin=1in}

\setcounter{section}{4}

\DeclareMathOperator{\range}{range}
\DeclareMathOperator{\sgn}{sgn}

% differential operator
\newcommand{\dd}{\mathop{}\!\mathrm{d}}

\newtheorem{thm}{Theorem}
\newenvironment{theorem}[1]{
	\renewcommand{\thethm}{#1}
	\begin{thm}
		}{\end{thm}}

\newtheorem{lma}{Lemma}
\newenvironment{lemma}[1]{
	\renewcommand{\thelma}{#1}
	\begin{lma}
		}{\end{lma}}

\theoremstyle{definition}
\newtheorem{ex}{}[section]
\newenvironment{exercise}[1]{
  \renewcommand{\theex}{\arabic{section}.\arabic{ex}}
	\begin{ex}
		}{\end{ex}}

\title{\S\arabic{section} Exercises}
\author{Connor Mooneyhan}
\date{May 25, 2025}

\begin{document}

\maketitle

\begin{ex}
	\textbf{A 1-form on $\mathbb{R}^3$} \\
	Let $\omega$ be the 1-form $z\dd x - \dd z$ and let $X$ be the vector field $y\partial / \partial x + x \partial / \partial y$ on $\mathbb{R}^3$.
	Compute $\omega(X)$ and $\dd\omega$.
\end{ex}
\begin{proof}
	Beginning with the formula as derived in the text, \begin{align*}
		\omega(X) & = \sum b_ia^i        \\
		          & = zy + 0 * x - 1 * 0 \\
		          & = zy.
	\end{align*}

	For the differential, \begin{align*}
		\dd\omega & = \dd(z) \wedge \dd x + \dd(-1)\dd z \\
		          & = \dd z \wedge \dd x + 0             \\
		          & = \dd z \wedge \dd x.
	\end{align*}
\end{proof}

\begin{ex}
	\textbf{A 2-form on $\mathbb{R}^3$} \\
	At each point $p \in \mathbb{R}^3$, define a bilinear function $\omega_p$ on $T_p(\mathbb{R}^3)$ by \begin{equation*}
		\omega_p(\mathbf{a}, \mathbf{b}) = \omega_p \left( \begin{bmatrix}a^1 \\ a^2 \\ a^3\end{bmatrix}, \begin{bmatrix}b^1 \\ b^2 \\ b^3\end{bmatrix} \right) = p^3 \det \begin{bmatrix}a^1 & b^1 \\ a^2 & b^2\end{bmatrix},
	\end{equation*}
	for tangent vectors $\mathbf{a}, \mathbf{b} \in T_p(\mathbb{R}^3)$, where $p^3$ is the third component of $p = (p^1, p^2, p^3)$.
	Since $\omega_p$ is an alternating bilinear function on $T_p(\mathbb{R}^3)$, $\omega$ is a 2-form on $\mathbb{R}^3$.
	Write $\omega$ in terms of the standard basis $\dd x^i \wedge \dd x^j$ at each point.
\end{ex}
\begin{proof}
	Since, \begin{align*}
		\omega_p(\mathbf{a}, \mathbf{b}) & = p^3 \det \begin{bmatrix} a^1 & b^1 \\ a^2 & b^2 \end{bmatrix} \\
		                                 & = p^3 (a^1b^2 - a^2b^1)
	\end{align*}
	and \begin{align*}
		(\dd x^1 \wedge \dd x^2)_p(\mathbf{a}, \mathbf{b}) & = (\dd x^1_p \wedge \dd x^2_p)(a, b)                  \\
		                                                   & = \dd x^1_p(a)\dd x^2_p(b) - \dd x^1_p(b)\dd x^2_p(a) \\
		                                                   & = a^1b^2 - b^1a^2,
	\end{align*}
	we have $\omega_p(\mathbf{a}, \mathbf{b}) = p^3(\dd x^1 \wedge \dd x^2)_p(\mathbf{a}, \mathbf{b})$.
\end{proof}

\begin{ex}
	\textbf{Exterior calculus} \\
	Suppose the standard coordinates on $\mathbb{R}^2$ are called $r$ and $\theta$ (this $\mathbb{R}^2$ is the $(r, \theta)$-plane, not the $(x, y)$-plane).
	If $x = r\cos\theta$ and $y=r\sin\theta$, calculate $\dd x, \dd y,$ and $\dd x \wedge \dd y$ in terms of $dr$ and $\dd\theta$.
\end{ex}
\begin{proof}
	\begin{align*}
		\dd x & = \dd(r\cos\theta)                                                                                    \\
		      & = \frac{\partial(r\cos\theta)}{\partial r}dr + \frac{\partial(r\cos\theta)}{\partial \theta}\dd\theta \\
		      & = \cos\theta dr - r\sin\theta \dd\theta,
	\end{align*}
	\begin{align*}
		\dd y & = \dd(r\sin\theta)                                                                                    \\
		      & = \frac{\partial(r\sin\theta)}{\partial r}dr + \frac{\partial(r\sin\theta)}{\partial \theta}\dd\theta \\
		      & = \sin\theta dr + r\cos\theta \dd\theta,
	\end{align*}
	\begin{align*}
		\dd x \wedge \dd y & = (\cos\theta dr - r\sin\theta \dd\theta) \wedge (\sin\theta dr + r \cos\theta \dd\theta)                                                                                          \\
		                   & = \cos\theta dr \wedge \sin\theta dr + \cos\theta dr \wedge r\cos\theta \dd\theta -r\sin\theta \dd\theta \wedge \sin\theta dr - r\sin\theta \dd\theta \wedge r\cos\theta \dd\theta \\
		                   & =  \cos\theta dr \wedge r\cos\theta \dd\theta -r\sin\theta \dd\theta \wedge \sin\theta dr                                                                                          \\
		                   & = r\cos^2\theta dr \wedge \dd\theta + r\sin^2\theta dr\wedge \dd\theta                                                                                                             \\
		                   & = r(\cos^2\theta + \sin^2\theta) dr\wedge \dd\theta                                                                                                                                \\
		                   & = r dr\wedge \dd\theta
	\end{align*}
\end{proof}

\begin{ex}
	\textbf{Exterior calculus} \\
	Suppose the standard coordinates on $\mathbb{R}^3$ are called $\rho$, $\phi$, and $\theta$.
	If $x = \rho\sin\phi\cos\theta$, $y = \rho\sin\phi\sin\theta$, and $z = \rho\cos\phi$, calculate $\dd x$, $\dd y$, $\dd z$, and $\dd x\wedge \dd y \wedge \dd z$ in terms of $\dd\rho$, $\dd\phi$, and $\dd\theta$.
\end{ex}
\begin{proof}
	\begin{align*}
		\dd x & = \dd(\rho\sin\phi\cos\theta)                                                                     \\
		      & = \sin\phi\cos\theta \dd\rho + \rho\cos\phi\cos\theta \dd\phi - \rho\sin\phi\sin\theta \dd\theta,
	\end{align*}
	\begin{align*}
		\dd y & = \dd(\rho\sin\phi\sin\theta)                                                                     \\
		      & = \sin\phi\sin\theta \dd\rho + \rho\cos\phi\sin\theta \dd\phi + \rho\sin\phi\cos\theta \dd\theta,
	\end{align*}
	\begin{align*}
		\dd z & = \dd(\rho\cos\phi)                        \\
		      & = \cos\phi \dd\rho - \rho\sin\phi \dd\phi,
	\end{align*}
	\begin{align*}
		\dd x \wedge \dd y \wedge \dd z = {} & - \sin\phi\cos\theta \dd\rho \wedge \rho\sin\phi\cos\theta \dd\theta \wedge \rho\sin\phi \dd\phi                \\
		                                     & + \rho\cos\phi\cos\theta \dd\phi \wedge \rho\sin\phi\cos\theta \dd\theta \wedge \cos\phi \dd\rho                \\
		                                     & + \rho\sin\phi\sin\theta \dd\theta \wedge \sin\phi\sin\theta \dd\rho \wedge \rho\sin\phi \dd\phi                \\
		                                     & - \rho\sin\phi\sin\theta \dd\theta \wedge \rho\cos\phi\sin\theta \dd\phi \wedge \cos\phi \dd\rho                \\
		= {}                                 & (\rho^2\sin^3\phi\cos^2\theta + \rho^2\cos^2\phi\cos^2\theta \sin\phi                                           \\
		                                     & + \rho^2\sin^3\phi\sin^2\theta + \rho^2\sin\phi\sin^2\theta \cos^2\phi) \dd\rho \wedge \dd\phi \wedge \dd\theta \\
		= {}                                 & \rho^2(\sin^3\phi\cos^2\theta + \cos^2\phi\cos^2\theta \sin\phi                                                 \\
		                                     & + \sin^3\phi\sin^2\theta + \sin\phi\sin^2\theta \cos^2\phi) \dd\rho \wedge \dd\phi \wedge \dd\theta             \\ =                        & \rho^2[\sin^3\phi(\cos^2\theta + \sin^2\theta) + \cos^2\phi\sin\phi (\cos^2\theta + \sin^2\theta)] \dd\rho \wedge \dd\phi \wedge \dd\theta \\
		= {}                                 & \rho^2(\sin^3\phi + \cos^2\phi\sin\phi)\dd\rho \wedge \dd\phi \wedge \dd\theta                                  \\
		= {}                                 & \rho^2\sin\phi(\sin^2\phi + \cos^2\phi)\dd\rho \wedge \dd\phi \wedge \dd\theta                                  \\
		= {}                                 & \rho^2\sin\phi \,\dd\rho \wedge \dd\phi \wedge \dd\theta                                                        \\
	\end{align*}
\end{proof}

\begin{ex}
	\textbf{Wedge product} \\
	Let $\alpha$ be a 1-form and $\beta$ a 2-form on $\mathbb{R}^3$. Then \begin{align*}
		 & \alpha = a_1 \dd x^1 + a_2 \dd x^2 + a_3 \dd x^3                                              \\
		 & \beta = b_1 \dd x^2 \wedge \dd x^3 + b_2 \dd x^3 \wedge \dd x^1 + b_3 \dd x^1 \wedge \dd x^2.
	\end{align*}
	Simplify the expression $\alpha \wedge \beta$ as much as possible.
\end{ex}
\begin{proof}
	\begin{align*}
		\alpha \wedge \beta & = a_1 \dd x^1 \wedge b_1 \dd x^2 \wedge \dd x^3 + a_2 \dd x^2 \wedge b_2 \dd x^3 \wedge \dd x^1 + a_3 \dd x^3 \wedge b_3 \dd x^1 \wedge \dd x^2 \\
		                    & = (a_1b_1 + a_2b_2 + a_3b_3)\dd x^1 \wedge \dd x^2 \wedge \dd x^3.
	\end{align*}
\end{proof}

\begin{ex}
	\textbf{Wedge product and cross product} \\
	The correspondence between differential forms and vector fields on an open subset of $\mathbb{R}^3$ in Subsection 4.6 also makes sense pointwise.
	Let $V$ be a vector space of dimension 3 with basis $e_1, e_2, e_3$, and dual basis $\alpha^1, \alpha^2, \alpha^3$.
	To a 1-covector $\alpha = a_1\alpha^1 + a_2\alpha^2 + a_3\alpha^3$ on $V$, we associate the vector $\mathbf{v}_\alpha = \left\langle a_1, a_2, a_3 \right\rangle \in \mathbb{R}^3$.
	To the 2-covector \begin{equation*}
		\gamma = c_1 \alpha^2 \wedge \alpha^3 + c_2 \alpha^3 \wedge \alpha^1 + c_3 \alpha^1 \wedge \alpha^2
	\end{equation*}
	on $V$, we associate the vector $\mathbf{v}_\gamma = \left\langle c_1, c_2, c_3 \right\rangle \in \mathbb{R}^3$.
	Show that under this correspondence, the wedge product of 1-covectors corresponds to the cross product of vectors in $\mathbb{R}^3$: if $\alpha = a_1 \alpha^1 + a_2 \alpha^2 + a_3\alpha^3$ and $\beta = b_1 \alpha^1 + b_2 \alpha^2 + b_3 \alpha^3$, then $\mathbf{v}_{\alpha \wedge \beta} = \mathbf{v}_\alpha \times \mathbf{v}_\beta$.
\end{ex}
\begin{proof}
	First, compute
	\begin{align*}
		\alpha \wedge \beta  = {} & a_1 b_2 \alpha^1 \wedge \alpha^2 + a_1 b_3 \alpha^1 \wedge \alpha^3 \\
		                          & +a_2b_1\alpha^2 \wedge \alpha^1 + a_2b_3 \alpha^2 \wedge \alpha^3   \\
		                          & +a_3b_1 \alpha^3 \wedge \alpha^1 + a_3 b_2 \alpha^3 \wedge \alpha^2 \\
		= {}                      & (a_2b_3 - a_3b_2)\alpha^2 \wedge \alpha^3                           \\
		                          & +(a_3 b_1 - a_1b_3)\alpha^3 \wedge \alpha^1                         \\
		                          & +(a_1b_2 - a_2b_1)\alpha^1 \wedge \alpha^2,
	\end{align*}
	so $\mathbf{v}_{\alpha \wedge \beta} = \left\langle a_2b_3 - a_3b_2, a_3b_1 - a_1b_3, a_1b_2 - a_2b_1 \right\rangle$, which is precisely the definition of $\left\langle a_1, a_2, a_3 \right\rangle \times \left\langle b_1, b_2, b_3 \right\rangle = \mathbf{v}_\alpha \times \mathbf{v}_\beta$.
\end{proof}

\begin{ex}
	\textbf{Commutator of derivations and antiderivations} \\
	Let $A = \oplus_{k=-\infty}^\infty A^k$ be a graded algebra over a field $K$ with $A^k = 0$ for $k < 0$.
	Let $m$ be an integer. A \textit{superderivation of $A$ of degree $m$} is a $K$-linear map $D: A \to A$ such that for all $k$, $D(A^k) \subset A^{k + m}$ and for all $a \in A^k$ and $b \in A^\ell$, \begin{equation*}
		D(ab) = D(a)b + (-1)^{km} a(Db).
	\end{equation*}
	If $D_1$ and $D_2$ are two superderivations of $A$ of respective degrees $m_1$ and $m_2$, define their \textit{commutator} to be \begin{equation*}
		[D_1, D_2] = D_1 \circ D_2 - (-1)^{m_1m_2}D_2 \circ D_1.
	\end{equation*}
	Show that $[D_1, D_2]$ is a superderivation of degree $m_1 + m_2$.
	(A superderivation is said to be \textit{even} or \textit{odd} depending on the parity of its degree.
	An even superderivation is a derivation; an odd superderivation is an antiderivation.)
\end{ex}
\begin{proof}
	Let $D_1$ and $D_2$ be superderivations of respective degrees $m_1$ and $m_2$.
	First, we show that $[D_1, D_2]$ is $K$-linear.
	Given $a, b \in A$, \begin{align*}
		[D_1, D_2](a + b) & = D_1 \circ D_2 (a + b) - (-1)^{m_1m_2}D_2 \circ D_1(a + b)                       \\
		                  & = D_1 (D_2(a) + D_2(b)) - (-1)^{m_1m_2}D_2(D_1(a) + D_1(b))                       \\
		                  & = D_1(D_2(a)) + D_1(D_2(b)) - (-1)^{m_1m_2}D_2(D_1(a)) - (-1)^{m_1m_2}D_2(D_1(b)) \\
		                  & = D_1(D_2(a)) - (-1)^{m_1m_2}D_2(D_1(a)) + D_1(D_2(b)) - (-1)^{m_1m_2}D_2(D_1(b)) \\
		                  & = [D_1, D_2](a) + [D_1, D_2](b),
	\end{align*}
	and for $r \in R$, \begin{align*}
		[D_1, D_2](ra) & = D_1(D_2(ra)) - (-1)^{m_1m_2}D_2(D_1(ra))  \\
		               & = rD_1(D_2(a)) - r(-1)^{m_1m_2}D_2(D_1(a))  \\
		               & = r(D_1(D_2(a)) - (-1)^{m_1m_2}D_2(D_1(a))) \\
		               & = r[D_1, D_2](a).
	\end{align*}
	Now suppose $a \in A^k$ and $b \in A^\ell$.
	First, observe that \begin{align*}
		[D_1, D_2](a) = {} & D_1(D_2(a)) - (-1)^{m_1m_2}D_2(D_1(a)),
	\end{align*}
	and $D_1(D_2(a)), D_2(D_1(a)) \in A^{k + m_1 + m_2}$. Thus, being a linear combination of these two, $[D_1, D_2](a) \in A^{k + m_1 + m_2}$.
	Finally, \begin{align*}
		[D_1, D_2](ab) = {} & D_1(D_2(ab)) - (-1)^{m_1m_2}D_2(D_1(ab))                                                         \\
		= {}                & D_1(D_2(a)b) + (-1)^{km_2}D_1(aD_2(b))                                                           \\
		                    & - (-1)^{m_1m_2}D_2(D_1(a)b) - (-1)^{m_1m_2 + km_1}D_2(aD_1(b))                                   \\
		= {}                & D_1 \circ D_2(a)b + (-1)^{km_1 + m_2m_1}D_2(a)D_1(b)                                             \\
		                    & + (-1)^{km_2}D_1(a)D_2(b) + (-1)^{km_2 + km_1}aD_1 \circ D_2(b)                                  \\
		                    & -(-1)^{m_1m_2}D_2 \circ D_1(a)b - (-1)^{m_1m_2 + km_2}D_1(a)D_2(b)                      \\
		                    & - (-1)^{m_1m_2 + km_1}D_2(a)D_1(b) - (-1)^{m_1m_2 + km_1 + km_2}aD_2 \circ D_1(b)                \\
		= {}                & \left[D_1 \circ D_2(a) - (-1)^{m_1m_2}D_2 \circ D_1(a)\right]b                                   \\
                        & + (-1)^{k(m_1 + m_2)}a\left[D_1 \circ D_2(b) - (-1)^{m_1m_2}D_2 \circ D_1(b)\right] \\
    = {}                & [D_1, D_2](a)b + (-1)^{k(m_1 + m_2)}a[D_1, D_2](b) \\
	\end{align*}
\end{proof}

\end{document}
